\section{Discussion}
\label{sec:discussion}

\subsection{Three Levels Tell Three Stories}

Our three experiments converge on a consistent but nuanced picture: compound concepts in \gpttwo are \emph{not} stored as dedicated directions, but they are also not \emph{merely} compositional.
The three analysis levels reveal complementary aspects of compound representation:

\begin{enumerate}[leftmargin=*,itemsep=0pt,topsep=0pt]
    \item \emph{Residual stream geometry} (cosine similarity 0.954): The compound representation is a subtle modification of the head noun direction. From the perspective of overall direction, ``washing machine'' looks like ``machine.''
    \item \emph{SAE feature decomposition} (56.3\% unique features): This subtle modification activates a largely different set of sparse features. The compound modifies the representation in ways that are invisible to cosine similarity but detectable by SAE decomposition.
    \item \emph{Next-token prediction} ($4{,}221\times$ priming): The model's primary mechanism for compound storage is sequential prediction. The concept of ``washing machine'' is encoded not as a direction but as a transition probability.
\end{enumerate}

This three-level picture resolves the apparent paradox: how can the compound be ``almost the same'' as the head noun (0.954 cosine similarity) yet activate mostly different SAE features (56.3\% unique)?
The answer is that cosine similarity measures the overall direction of a 768-dimensional vector, while SAE features capture fine-grained structure along specific feature directions.
A small rotation in the overall direction can correspond to large changes in which features are active, because SAE features operate near their activation thresholds.

\subsection{The Role of Compositionality}

The correlation between linguistic compositionality and unique SAE features (Cohen's $d = 2.54$) reveals that the model allocates more specialized representation to semantically opaque compounds.
``Red herring'' activates 82.1\% unique features because the idiomatic meaning (a misleading clue) cannot be derived from ``red'' and ``herring'' separately---the model must encode this non-compositional semantics through additional features.
In contrast, ``kitchen chair'' (46.9\% unique) is nearly derivable from its parts, requiring less specialized representation.

An intriguing pattern emerges from the priming data: ``guinea pig'' has the highest priming ratio ($28{,}037\times$) but moderate unique features (61.7\%), while ``red herring'' has the most unique features (82.1\%) but minimal priming ($4.3\times$).
This dissociation reflects two independent dimensions of compound representation: statistical predictability (how frozen the compound is in text) and semantic opacity (how much the compound meaning deviates from component meanings).
``Guinea pig'' is statistically frozen but semantically processed compositionally by the SAE, while ``red herring'' is statistically free but semantically specialized.

\subsection{Implications}

\para{For mechanistic interpretability.}
Our findings challenge the assumption that cosine similarity is sufficient for detecting concept-specific processing.
Two representations with 0.895 cosine similarity (``red herring'' vs.\ ``herring'') can have only 14.3\% feature overlap, suggesting that interpretability methods relying solely on geometric measures may miss important structure.
SAE decomposition provides a complementary and more sensitive lens, though one must account for potential SAE artifacts~\citep{chanin2024absorption}.

\para{For the superposition hypothesis.}
Compound concepts illustrate how LLMs manage the ``too many concepts, too few dimensions'' problem.
Rather than allocating dedicated directions for each of the millions of possible compounds, the model constructs multi-token concepts sequentially through transition probabilities.
This is a fundamentally different strategy from superposition: instead of packing more features into the same space, the model distributes concept representation across \emph{time steps}.
Superposition handles the single-token feature bottleneck; sequential construction handles the multi-token concept bottleneck.

\para{For concept editing.}
Editing compound concept knowledge (\eg removing the model's knowledge that a washing machine is an appliance) would likely require intervention at multiple levels: modifying the transition probabilities (attention patterns and FFN weights that produce the $4{,}221\times$ priming effect) \emph{and} the compound-specific SAE features (the 57.9\% unique features for ``washing machine'').
Modifying a single direction is unlikely to suffice, given that compound representations are distributed across sequential prediction and fine-grained feature modulation.

\subsection{Limitations}

\para{Single model.}
We study \gpttwo (124M parameters) only.
Larger models with more dimensions may allocate dedicated directions for common compounds.
The absence of holistic compound representations in GPT-2 may reflect capacity constraints rather than a general principle.

\para{Minimal context.}
We use the context ``The [compound]'' for all experiments.
Richer contexts (\eg ``She loaded the washing machine with clothes'') could produce different activation patterns.
Natural corpus contexts from large-scale text would better capture the range of compound usage.

\para{SAE artifacts.}
Pre-trained SAEs have known issues including feature absorption~\citep{chanin2024absorption} and polysemantic features.
The 56.3\% unique feature fraction might partly reflect SAE reconstruction artifacts rather than genuine compound-specific computation.
However, the strong correlation between unique features and compositionality ($d = 2.54$) suggests that the pattern is meaningful even if exact percentages are noisy.

\para{Hand-assigned compositionality ratings.}
Our compositionality labels are hand-assigned rather than derived from human judgment experiments.
Validated compositionality ratings from psycholinguistic studies~\citep{reddy2011empirical} would strengthen the compositionality analysis.

\para{Tokenization effects.}
Some compounds split unevenly under BPE tokenization (\eg ``coffee'' $\to$ \texttt{co} + \texttt{ffee}), meaning the ``modifier'' representation already involves compositional processing.
This complicates the interpretation of feature overlap statistics for multi-token modifiers.
